\documentclass[11pt,a4paper]{article}

\usepackage[utf8]{inputenc}
\usepackage[T1]{fontenc}
\usepackage{lmodern}
\usepackage[margin=2.5cm]{geometry}
\usepackage{parskip}
\usepackage{ragged2e}
\usepackage[hidelinks]{hyperref}

\pagestyle{empty}
\setlength{\parindent}{0pt}

\begin{document}

% Candidate information
{\Large\bfseries John Doe}\\[0.3cm]


\href{mailto:john@example.com}{john@example.com}\\



+33 6 00 00 00 00\\



Paris\\


\vspace{1cm}

% Recipient
\begin{flushright}

\textbf{Clarins}\\



Pontoise\\


\vspace{0.5cm}

12 novembre 2026

\end{flushright}

\vspace{0.8cm}

% Subject
\textbf{Subject: Candidature spontanée – Stage en Data Engineering}

\vspace{0.6cm}

% Salutation

Madame, Monsieur,


\vspace{0.5cm}

% Cover letter body
\justifying


Je vous adresse ma candidature spontanée pour un stage en Data Engineering au sein de Clarins à Pontoise. Étudiant en Master informatique (2024–2026), je m’intéresse particulièrement à l’ingénierie des données et à sa contribution à l’innovation, en résonance avec votre intérêt pour la data analytics dans le secteur des cosmétiques.

\par

Lors d’un précédent stage en développement, j’ai réalisé des scripts en Python et écrit des requêtes SQL, dans un environnement Linux. Ces expériences me donnent de solides bases pour contribuer à des tâches de collecte et de transformation de données, d’automatisation par scripts et de requêtage fiable, tout en portant attention à la qualité et à la clarté du code.

\par

Motivé et curieux, je souhaite mettre ces compétences au service de vos équipes et apprendre au contact de vos pratiques. Je serais ravi d’échanger pour préciser vos besoins et la manière dont je peux y contribuer.

\par


\vspace{0.7cm}

% Closing

Cordialement,


\vspace{0.5cm}

\textbf{John Doe}

\end{document}

